% ============================================================
% Fusion Compiler 第 2 章相对 ICC2 新增/显著不同的节（可 \input 片段）
% 原书：Preparing the Design（约 p.80–103、p.135–137）
% ============================================================

\section{指定库子集限制}
\subsection*{Specifying Library Subset Restrictions}

可将时序单元（sequential cells）以及已例化组合单元（instantiated combinational cells）的映射与优化限制到参考库与库单元的某个子集。要指定一个或多个库子集限制，先使用 \cmd{define\_libcell\_subset}，再使用 \cmd{set\_libcell\_subset}。

\begin{noteBox}
库子集限制适用于叶子单元（leaf cells），不适用于层次单元（hierarchical cells）。
\end{noteBox}

要在叶子单元上设置库子集限制：
\begin{enumerate}
  \item 用 \cmd{define\_libcell\_subset} 定义库单元子集。\\
    库单元被分入子集后，不再可用于优化期间的一般映射，而仅用于时序单元与已例化组合单元。
  \item 用 \cmd{set\_libcell\_subset} 限制时序单元与已例化组合单元的优化。\\
    优化期间，工具使用步骤 1 中定义的库子集来优化时序单元（已映射与未映射）以及已映射的组合单元。
\end{enumerate}

下例中，\cmd{define\_libcell\_subset} 将 \texttt{SDFLOP1} 与 \texttt{SDFLOP2} 库单元归入名为 \texttt{special\_flops} 的子集，随后 \cmd{set\_libcell\_subset} 将叶子单元 \texttt{LEAF1} 的映射限制到该 \texttt{special\_flops} 库子集：
\begin{lstlisting}
fc_shell> define_libcell_subset \
   -libcell_list "SDFLOP1 SDFLOP2" -family_name special_flops
fc_shell> set_libcell_subset \
   -object_list "HIER1/LEAF1" -family_name special_flops
\end{lstlisting}

\subsubsection{报告库子集限制}
\subsubsection*{Reporting Library Subset Restrictions}

要报告为时序单元与已例化组合单元指定的库子集，使用 \cmd{report\_libcell\_subset}。

\subsubsection{移除库子集限制}
\subsubsection*{Removing Library Subset Restrictions}

要移除库子集，使用 \cmd{remove\_libcell\_subset}。

% ============================================================
